\chapter{The Memory Model}
\label{ch:memory-model}

\epigraph{``All problems in computer science can be solved by another level of indirection, except for the problem of too many levels of indirection.''}{David Wheeler}

\learninggoal{
	\item Describe the virtual memory abstraction and its components: pages, frames, and page tables.
	\item Distinguish between stack and heap memory organisation and their lifetime semantics.
	\item Explain how memory permissions (read, write, execute) are enforced at the hardware level.
	\item Understand segmentation, paging, and the address translation pipeline.
	\item Relate memory model features to security properties and vulnerabilities.
}

\section{Virtual Memory}

Every modern processor provides a \emph{virtual memory} abstraction. Each process sees a private, contiguous address space starting at address~0. The processor and operating system cooperate to map these virtual addresses to physical memory locations.

\begin{definition}[Virtual Memory]
	\textbf{Virtual memory} is an indirection layer that maps virtual addresses (used by programs) to physical addresses (used by hardware). The mapping is performed by the \emph{memory management unit} (MMU) using data structures called \emph{page tables}.
\end{definition}

Why virtual memory? Three reasons:

\begin{enumerate}
	\item \textbf{Isolation:} process A cannot read or write process B's memory, because A's virtual addresses map to different physical pages.
	\item \textbf{Security:} the kernel can mark pages as non-executable, non-writable, or non-readable. Any violation triggers a processor exception.
	\item \textbf{Convenience:} each process can use a flat address space without coordinating with other processes or worrying about physical memory fragmentation.
\end{enumerate}

\begin{principle}[The MMU as Security Enforcer]
	Every memory access goes through the MMU. This means that memory protections are not optional---they are enforced by the processor on \emph{every} load and store instruction, including those in kernel code. A bug in user space cannot bypass the MMU (barring processor vulnerabilities like Meltdown).
\end{principle}

\section{Paging}

The dominant virtual memory implementation is \emph{paging}. Virtual and physical memory are divided into fixed-size blocks called \emph{pages} (typically 4\,KB on x86-64).

\begin{definition}[Page, Frame, Page Table]
	A \textbf{page} is a fixed-size block of virtual memory. A \textbf{frame} is a fixed-size block of physical memory of the same size. A \textbf{page table} is a multi-level data structure that maps virtual pages to physical frames.
\end{definition}

x86-64 uses four levels of page tables (PML4, PDPT, PD, PT) for 4\,KB pages. A virtual address is split into fields that index each level:

\begin{lstlisting}[caption=Virtual address layout on x86-64 (4KB pages)]
  Virtual address (48 bits):
  | PML4 (9) | PDPT (9) | PD (9) | PT (9) | Offset (12) |
\end{lstlisting}

Each level's entry contains the physical address of the next-level table, plus permission bits:

\begin{itemize}
	\item \texttt{P} (Present): the page is in physical memory.
	\item \texttt{R/W}: writable (0 = read-only).
	\item \texttt{U/S}: user-accessible (0 = supervisor-only).
	\item \texttt{NX} (No-Execute): code cannot be executed from this page.
	\item \texttt{A} (Accessed): set by hardware on first access.
	\item \texttt{D} (Dirty): set by hardware on first write.
\end{itemize}

\begin{example}[Address Translation]
	Consider a load from virtual address \texttt{0x7fff00123456}. The processor:
	\begin{enumerate}
		\item Extracts bits 47:39 to index PML4 (entry \texttt{0x1FF}).
		\item Reads the PML4 entry, which contains the physical address of the PDPT.
		\item Extracts bits 38:30 to index the PDPT.
		\item Repeats for PD (bits 29:21) and PT (bits 20:12).
		\item Reads the PTE, which contains the physical frame number.
		\item Adds the 12-bit page offset (bits 11:0) to form the final physical address.
	\end{enumerate}
	If any page-table entry has \texttt{P}=0, the processor raises a page fault. If \texttt{U/S}=0 while the processor is in user mode, a protection fault occurs.
\end{example}

\section{Memory Layout of a Process}

A typical process on Linux has the following virtual memory layout:

\begin{lstlisting}[caption=Linux process memory layout (x86-64)]
  High addresses (0x7fffffffffff)
  +---------------------------+
  | Kernel space              |  Not accessible from user mode
  | (mapped in every process) |
  +---------------------------+
  | Stack                     |  Grows downward (toward lower addresses)
  | (per-thread, ~8MB by def)|
  +---------------------------+
  | Memory mapping region     |  mmap, shared libraries, heap expansion
  | (grows downward)          |
  +---------------------------+
  | Heap                      |  Grows upward (brk/sbrk)
  +---------------------------+
  | BSS                       |  Uninitialized global data
  +---------------------------+
  | Data                      |  Initialized global data
  +---------------------------+
  | Text (code)               |  Executable machine code
  | (read-only, executable)   |
  +---------------------------+
  Low addresses (0x400000)
\end{lstlisting}

\subsection{The Stack}

The stack is a contiguous region of memory that grows downward. It stores:

\begin{itemize}
	\item \textbf{Return addresses:} where to resume execution after a function call.
	\item \textbf{Local variables:} automatic storage duration.
	\item \textbf{Saved frame pointers:} the previous base pointer (\texttt{rbp}) value.
	\item \textbf{Function arguments:} passed in registers or on the stack.
\end{itemize}

\begin{definition}[Stack Frame]
	A \textbf{stack frame} is the region of the stack allocated for a single function invocation. It contains the function's local variables, saved registers, and the return address. The exact layout is determined by the calling convention (Section~\ref{sec:calling-conventions}).
\end{definition}

The stack pointer (\texttt{rsp}) points to the top of the stack. The base pointer (\texttt{rbp}) typically points to the start of the current frame. On x86-64, arguments are passed in registers for the first six integer arguments (\texttt{rdi}, \texttt{rsi}, \texttt{rdx}, \texttt{rcx}, \texttt{r8}, \texttt{r9}); additional arguments go on the stack.

\subsection{The Heap}

The heap is a dynamically allocated memory region. Unlike the stack, heap allocations:

\begin{itemize}
	\item Have arbitrary lifetimes (must be explicitly freed).
	\item Can be resized with \texttt{realloc}.
	\item Are managed by an allocator (e.g., \texttt{ptmalloc} on Linux).
	\item Are subject to fragmentation.
\end{itemize}

\begin{definition}[Heap Allocator]
	A \textbf{heap allocator} manages the heap region, satisfying \texttt{malloc} and \texttt{free} requests. It maintains internal metadata (chunk headers, free lists) to track which regions are in use. Corrupting this metadata is a common exploitation technique.
\end{definition}

\section{Memory Protection and the NX Bit}

Before the NX (No-Execute) bit was introduced, all memory pages were implicitly executable. This meant that an attacker who could write data to memory (via a buffer overflow) could then jump to that data and execute it as code. The NX bit, introduced in x86-64, allows each page to be independently marked executable or non-executable.

\begin{principle}[W\textasciicircum X (Write XOR Execute)]
	A memory page should be either writable or executable, but never both. This principle, called \textbf{W\textasciicircum X}, prevents the most direct exploitation path: writing shellcode to memory and executing it.
\end{principle}

The stack and heap are mapped as readable and writable, but not executable. The code section (\texttt{.text}) is mapped as readable and executable, but not writable. An attacker cannot write new code to memory and execute it---they must reuse existing code (Section~\ref{ch:exploitation}).

\section{TLB and Caching}

Virtual-to-physical address translation would be slow if every memory access required four page-table walks. The \emph{translation lookaside buffer} (TLB) caches recent translations.

\begin{definition}[TLB]
	The \textbf{TLB} is a hardware cache in the CPU that stores recent virtual-to-physical page mappings. A TLB hit avoids the page-table walk; a TLB miss triggers a hardware walk (or a microcode-assisted walk).
\end{definition}

TLBs are small (typically 32--64 entries for L1, up to 2048 for L2). Context switches (switching between processes) flush the TLB, which is why system calls and thread switches have a performance cost.

\section{Security Implications of the Memory Model}

Every feature of the memory model has security implications:

\begin{longtable}{p{3.5cm}p{4.5cm}p{4.5cm}}
	\toprule
	\textbf{Feature}   & \textbf{Security benefit}                        & \textbf{Security risk}                            \\
	\midrule
	Virtual memory     & Process isolation prevents cross-process attacks & Information leaks via shared pages (e.g., KASLR)  \\
	Page permissions   & NX prevents direct shellcode execution           & WX pages in buggy drivers bypass NX               \\
	TLB                & Fast address translation                         & TLB side channels can leak address information    \\
	Heap allocator     & Transparent memory management                    & Heap metadata corruption leads to arbitrary write \\
	Stack frame layout & Efficient function call mechanism                & Return address overwrite hijacks control flow     \\
	\bottomrule
\end{longtable}

\section{Exercises}

\begin{exercise}
	On x86-64 with 4\,KB pages, how many bytes are used for the virtual address in each of the four page-table fields (PML4, PDPT, PD, PT)? What fraction of the 64-bit address space is currently usable?
\end{exercise}

\begin{exercise}
	Given a virtual address \texttt{0x7f0000000000}, compute the indices for PML4, PDPT, PD, and PT. Assume 4\,KB pages and 48-bit virtual addresses.
\end{exercise}

\begin{exercise}
	Why do large pages (2\,MB or 1\,GB) reduce TLB pressure? What is the trade-off for security?
\end{exercise}

\begin{exercise}
	Research the KASLR (Kernel Address Space Layout Randomisation) feature. How does it randomise the kernel's position in virtual memory? How have attackers bypassed it using TLB side channels?
\end{exercise}
